\documentclass[a4paper,11pt]{article}% {{{
% -------------------------------------------------
% Set up the page margins.
% -------------------------------------------------
\usepackage[margin=1in]{geometry}
% \usepackage[headheight = .5in, headsep = \baselineskip, top = 1in, left = 1in, right = 1in, textwidth = 7.52 in]{geometry}
% \usepackage[text={6.5in,9.5in},centering]{geometry}
% \usepackage[margin=1in]{geometry}
% \setlength{\topmargin}{-0.25in}

% -------------------------------------------------
% Load Packages here
% -------------------------------------------------
\usepackage{amssymb, latexsym, amsmath, graphics, fullpage, epsfig, amsthm, relsize, pgf, tikz, amsfonts, makeidx, latexsym, ifthen, calc}
\usepackage[colorlinks,citecolor=blue,urlcolor=blue]{hyperref}
%\usepackage{eucal} this is a different font than \mathcal{•}
\usetikzlibrary{arrows}
\usepackage{mathrsfs}
\usepackage[shortlabels]{enumitem}
\usepackage{tikz}
\usepackage{amsmath}
\usetikzlibrary{arrows}
\usetikzlibrary{shadows.blur}
\usepackage{pgfplots}
\usepackage{mwe} % For dummy images
\usepackage{subcaption}
\usepackage{multirow}
\usepackage{dcolumn}
\newcolumntype{2}{D{.}{}{2.0}}
\usepackage{multicol}
\numberwithin{equation}{section}

% -------------------------------------------------
% check references
% -------------------------------------------------
% \usepackage{refcheck}
\usepackage[notref,notcite]{showkeys}
% \usepackage[notcite]{showkeys}
% \usepackage{showkeys}


% -------------------------------------------------
% Environments here
% -------------------------------------------------
\theoremstyle{plain}
\newtheorem{theorem}{Theorem}[section]
\newtheorem{corollary}[theorem]{Corollary}
\newtheorem{lemma}[theorem]{Lemma}
\newtheorem{proposition}[theorem]{Proposition}

\theoremstyle{definition}
\newtheorem{definition}[theorem]{Definition}
\newtheorem{hypothesis}[theorem]{Hypothesis}
\newtheorem{assumptions}[theorem]{Assumption}

\newtheorem{remark}[theorem]{Remark}
\newtheorem{conjecture}[theorem]{Conjecture}
\newtheorem{example}[theorem]{Example}
\newtheorem{notation}[theorem]{Notation}


% -------------------------------------------------
%  Define short hand symbols.
% -------------------------------------------------

\newcommand{\lip}{\textup{Lip}_b}
\newcommand{\liph}{\text{Lip}_{\hat\rho}}
\newcommand{\R}{\mathbb{R}}
\newcommand{\E}{\mathbb{E}}
\newcommand{\e}{\mathrm{e}}
\newcommand{\Norm}[1]{\left\|  #1   \right\|}
\newcommand{\rhoNorm}[1]{\left|  #1   \right|}
\newcommand{\ud}{\ensuremath{ \mathrm{d} }}


\title{Comparison to previous results}
\allowdisplaybreaks
%}}}

\begin{document}
I will be making comparisons of our paper to the following four papers: \cite{TessitoreZabczyk98M} (1998), \cite{AM03} (2003), \cite{MSY15} (2015) and \cite{MSY16} (2016). First I will summarize what is shown in each of the above papers.

\tableofcontents

\section{Summary of our paper.}

Our paper handles the SPDE
	\begin{equation}\label{E:gspde}
		\begin{cases}
			\frac{\partial u}{\partial t}(t,x) - \frac{1}{2}\Delta u(t,x) = b(x,u(t,x))\dot{W}(t,x) & x\in \R^d, \; t>0
			\\ u(0,\cdot) = \mu(\cdot)
		\end{cases}.
	\end{equation}
The noise $W$ is a zero mean Gaussian processes defined on the space of test functions $C^{\infty}_0([0,\infty) \times \R^d)$ that is white in time and colored in space with covariance function
	\[
		J(\varphi, \psi) = \int_0^\infty \ud s \int_{\R^d} \ud x \int_{\R^d} \ud y \; \psi(s,y) f(x-y) \varphi(x-y).
	\]
The function $f$ is referred to as the spatial correlation function and its Fourier transform $\mathcal{F}f = \hat{f}$ is referred to as the spectral density.

	\begin{remark}
		At the moment, I am unsure if we can use a space-time white noise since the moment bounds proven in \cite{CH16Comparison} are only for a noise that is white in time and colored in space.
		\textcolor{magenta}{Our results do cover space-time white noise case when $d=1$. The corresponding moment formula is
		much more explicit;  see \cite{ChenDalang13Heat}.}
	\end{remark}

	\begin{remark}
		The following approach has been used to prove the existence of an invariant measure. First we show that the semigroup generated by $\Delta$ satisfies the Feller property. Second, we show that the solution is bounded in probability in $L^2_{\rho}(\R^d)$.

		The first has been shown to be true for a wide range of operators which include $\Delta$. The second step is the step that requires some work and the trick is often to show that
			\[
				\sup_{t \ge 0} \mathbb{E}\Norm{u(t\cdot)}_{L^2_{\rho}}^2 < \infty
			\]
		and, in fact, this is what is done in \cite{AM03, MSY15,MSY16,TessitoreZabczyk98M}.
	\end{remark}

The following assumptions must be made.
	\begin{assumptions}
		\label{A:General}
		\begin{enumerate}[(i)]
			\item The initial measure $\mu$ satisfies the following integrability condition:
			\begin{equation}
			\int_{\R^d}\exp\left(-a|x|^2\right)|\mu|(\ud x) < \infty \quad \text{for any $a>0$;} \label{A:icon}
			\end{equation}
			\item The spectral measure $\hat{f}$ satisfies {\it Dalang's condition}:
			\begin{equation}
			\Upsilon(\beta) := (2\pi)^{-d}\int_{\R^d}\dfrac{\hat f(\ud \xi)}{\beta + |\xi|^2} <
			\infty \quad \text{for some and therefore any $\beta>0$}  \label{A:dalcon}
			\end{equation}
			\item The diffusion coefficient $b(x,u)$ satisfies the following Lipschitz and boundedness conditions:
			\begin{equation}
			|b(x,u)-b(x,v)|< L |u-v| \quad \text{and} \quad |b(x,0)| \le L_0 \quad u,v \in \R, \; x \in \R^d \label{A:lipcon}.
			\end{equation}
		\end{enumerate}
	\end{assumptions}

Our first main result handles the case of the existence of an invariant measure for $L^2_{\rho}(\R^d)$ initial conditions, namely the following SPDE

\begin{equation}
\begin{cases}
\dfrac{\partial u}{\partial t}(t,x)-\frac{1}{2}\Delta_x u(t,x) = b(x,u(t,x)\dot W(t,x) & \text{$x\in \R^d$ , $t>0$}
\\ u(0,x)=\varphi(x) & \varphi(x) \in L^2_\rho(\R^d)
\end{cases}  \label{E:spde}.
\end{equation}

Before we state the result, we need to make an additional assumption on the spectral measure.

\begin{assumptions}
	\item \begin{equation}
	\int_0^s \ud r \: (s-r)^{-2\alpha} \int_{\R^{d}} \hat f(\ud \xi) \exp\left(-(s-r)|\xi|^2 \right) < \infty \label{A:fcon}
	\end{equation}
\end{assumptions}

With this assumption now stated, we now state our first main result, Theorem \ref{T:Main}.

\begin{theorem}
	\label{T:Main}
	Assume that Assumption \ref{A:General} holds with the stronger assumption that $\Upsilon(0) <\infty$. Also assume \eqref{A:fcon}. Now for any admissible weight $\rho$, suppose there exists another admissible weight $\hat\rho$ and an element $\hat{\varphi} \in L^2_\rho(\R^d) \cap L^2_{\hat\rho}(\R^d)$ such that
	\begin{align*}
	\int_{\R^d}\dfrac{\rho(x)}{\hat\rho(x)}\ud x < \infty \quad \text{and $u^{\hat\varphi}(\cdot,x)$ is bounded in probability in $L^2_{\hat\rho}(\R^d)$.}
	\end{align*}
	Then for any $t_0 >0$ there exists an invariant measure for the laws of the solution of the SPDE \eqref{E:gspde}, $\{\mathscr L (u^{\hat\varphi}(t,\cdot))\}_{t \ge t_0}$,  in $L^2_\rho(\R^d)$ for any $L^2_{\rho}(\R^d)$ initial condition.
\end{theorem}

We now state our last assumption which will be required for the proof of one of our last main result, Theorem \eqref{T:Main2}.

\begin{assumptions} $\;$
	\begin{equation}
	\sup_{x\in \R^d}	\big(G(t,\cdot)*|\mu|)(x) \big) < K(t), \quad \forall t>0 \text{ and } K(t)\in \mathbb N \label{A:newgicon}
	\end{equation}
\end{assumptions}

\begin{remark}
	It is easy to see that Assumption \eqref{A:newgicon} is satisfied whenever $\mu=\varphi \in L^2_\rho(\R^d)$ is bounded and also when $\mu = \delta_0$ where $\delta_0$ is the Dirac delta measure.
\end{remark}

\begin{theorem}
	\label{T:Moment}
	Suppose Assumption \ref{A:General} holds.
	Let $u(t,x)$ be the unique solution to \eqref{E:spde} starting from  $\varphi(\cdot) \in L^2_\rho$. Then for all $T>0$ and for all $t \in [0,T]$, $u(t,\cdot)\in L^2_\rho(\R^d)$ with
	\begin{align}
	\label{E:uNormRho}
	\E\left(| u(t,\cdot)|_\rho^2\right)
	\le  2 C_T \left[ \Norm{\rho}_{L^1(\R^d)} + C_\rho(T) |\varphi|_{\rho}^2 \right]<\infty.
	\end{align}
	where we use $C_T:= H(T,\gamma_2)$. Moreover, if $\Upsilon(0)<\infty$ and assume \eqref{A:newgicon}, in other words, for each $t>0$ there exists a $K(t) \in \mathbb N$ such that $\sup\limits_{x\in \R^d}\big(|\varphi|*G(t,\cdot)\big)(x) < K(t)$, then one can show
	\begin{align*}
	\sup_{t\ge 0}\E\left(|u(t,\cdot)|_\rho^2\right) < \infty
	\end{align*}
\end{theorem}

The following Theorem proves the existence of an invariant measure for measure valued initial conditions.

\begin{theorem}
	\label{T:Main2}
	Suppose the assumptions of Theorem \ref{T:Main} remain true and in addition assume \eqref{A:newgicon} holds true. Suppose there existed
	another admissible weight $\hat{\rho}$ such that
	\begin{align*}
	\int_{\R^d}\dfrac{\rho(x)}{\hat\rho(x)}\: \ud x < \infty.
	\end{align*}
	Then there exists an invariant measure for the laws of the solutions to
	\eqref{E:gspde} in $L^2_\rho(\R^d)$ for $t>l>0$. In particular, for any $l>0$ there exists a measure that is
	invariant for $\{\mathscr L (u^\mu(t,\cdot))\}_{t>l}$ where $u$ is the solution to \eqref{E:gspde} with
	initial measure $\mu$.
\end{theorem}

\section{Summaries of previous papers.}

	\subsection{Invariant Measures for Stochastic Heat Equations. Tessitore and Zabczyk \cite{TessitoreZabczyk98M}}
This paper works with the following SPDE:
	\[
		\begin{cases}
		\partial_t X(t,\xi) = \frac{1}{2} \Delta X(t,\xi) + b(\xi,X(t,\xi)) \dot{W}(t,\xi), & t>0, \; \xi \in \R^d
		\\ X(0,\xi) = x(\xi) & x \in \R^d
		\end{cases}
	\]
Here $W$ is a spatially homogeneous Wiener processes on $(\Omega,\mathscr{F},(\mathscr{F}_t), \mathbb{P})$. It has a unbounded covariance operator $Q$ of convolution type in the sense that
	\[
		Qx = \Gamma*x
	\]
where $\Gamma$ is a positive definite distribution with even spectal density $\gamma \ge 0$. We have
	\[
		\mathbb{E}\langle W(t),\phi \rangle\langle W(s) , \psi \rangle = \min(t,x) \langle \Gamma * \phi , \psi \rangle
	\]
for Schwarz functions $\phi$ and $\psi$.

The following theorem shows the existence of an invariant measure when assuming a number of assumptions.

	\begin{theorem}[Theorem 3.1:]\label{T:TZth3.1}
Assume that:
		\begin{enumerate}
			\item There exists a $p\in [1,\infty],$ with $(1-1/p)(d/2) < 1$, such that $\gamma \in L^p(\R^d)$.
			\item $|b(x,y)-b(x,z)| \le l|y-z| \quad \text{and} \quad |b(x,0)| \le l_0 \quad \forall \; x \in \R^d, \; y,z \in \R.$
		\end{enumerate}
In addition, assume that for a given admissible weight, $\rho$, there exists another admissible weight, $\hat{\rho}$ such that $\rho \hat{\rho}^{-1} \in L^1(\R^d)$ and that $X_{\hat{\rho}}(t,\cdot)$ is bounded in probability in $L^2_{\hat{\rho}}$. Then there exists an invariant measure for the solution to the SPDE in question in $L^2_{\rho}$.
	\end{theorem}

The following and last theorem from this paper proves that for the constant initial condition, $x(\xi)=1$, that the solution to the SPDE is bounded in probability.

	\begin{theorem}[Theorem 3.3]\label{T:TXth3.3}
Assume $(a)$ and $(b)$ from Theorem \ref{TZth3.1} and define
		\[
			\tilde{\Gamma}|\widehat{\gamma^{1/2}} | * |\widehat{\gamma^{1/2}}| \quad \text{and} \quad \textbf{l}= \frac{\mathbf{\Gamma}(d/2-1) } {4\pi^{d/2}} \int_{\R^d} \tilde{\Gamma}(\zeta) |\zeta|^{2-d} \ud \zeta,
		\]
where $\mathbf{\Gamma}$ is the gamma function. If $d \ge 3$ and $\textbf{l} < l^{-2}$, where $l$ is the Lipschitz constant of $b$. Then for all admissible weights $\rho$,
		\[
			\sup_{t\ge0} \mathbb{E}(|X_1(t,\cdot)|_{\rho}^2) < \infty
		\]
	\end{theorem}

\begin{remark}
	Theorem \ref{T:TZth3.1} can only be explicitly applied to the SPDE with constant one initial condition since the solution has only been shown to be bounded in probability for this case which is the content of Theorem \ref{T:TXth3.3}.
\end{remark}

\subsection{Invariant measures for the stoch. heat equations with unbounded coefficients. Assing and Manthey \cite{AM03}}

This paper works with the following spde:
	\[
	\begin{cases}
		\frac{\partial}{\partial t} u(t,x) = \lambda \Delta u(t,x) + f(u(t,x)) + \sigma(u(t,x)) \cdot \dot{W}(t,x), & t>0, \; x \in \R^d
	\\ u(0,x) = \theta(x) & x \in \R^d
	\end{cases}
	\]
where $\lambda >0$ and $\dot{W}$ is white in time and either white or colored in space ($Q$-Wiener Processes). The noise is defined as follows: consider a sequence of $\mathbb{F}$ adapted independent Wiener processes, $(B_k(t))_k$, on a complete probability space $(\Omega, \mathscr{F}, \mathbb{P})$, and $\mathbb{F} = (\mathscr{F}_t)_{t\ge 0}$ denotes a right-continuous filtration of $\mathscr{F}$ and $\mathscr{F}_0$ contains all sets of $\mathbb{P}$ measure 0 and $B_k(t)-B_s(t)$ is independent of $\mathscr{F}_s$ with $k>s$. Now fix an orthonormal basis, $(e_k(x))_k)$, of $L^2(\R^d)$ that is uniformly bounded in the sense that
	\[
		\sup_{k} \sup_{x\in \R^d} \Norm{ e_k(\cdot) }_{L^\infty(\R^d)} < \infty.
	\]
There are two cases when defining the noise.

\begin{enumerate}
	\item \textbf{Cylindrical case $(d=1)$:}
		\[
			W(t) = \sum_{k=1}^{\infty} B_k(t)e_k
		\]
	\item \textbf{Nuclear case $(d \ge \textcolor{magenta}{1})$:}
		\[
			W(t) = \sum_{k=1}^{\infty} \sqrt{a_k} B_k(t) e_k
		\]
	where $(a_k)_k$ is a sequence of nonnegative real numbers satisfying $\sum a_k < \infty$. In the Nuclear case we have $a := \sum a_k \Norm{e_k}_{L^\infty(\R^d)} < \infty$.
\end{enumerate}

\bigskip
\textcolor{magenta}{Here we have a main contribution that our noise go beyond nuclear case for any dimension. Riesz
kernel is not even a trace-class operator. We should emphasize this contribution.}
\bigskip

In the paper \textcolor{magenta}{they} only consider the admissible weight
	\[
		\rho(x) = (1+|x|^2)^{-\varrho/2} \quad \varrho > d.
	\]
Here is the first result from this paper.
	\begin{theorem}[Theorem 1:]\label{T:AMth1}
Assume that
		\begin{enumerate}
			\item $|\sigma(x) - \sigma(y)| \le c_\sigma |x-y|$
			\item $ |f(x)-f(y)| \le c_{\nu}|x-y|(1+|x|^{\nu -1} + |y|^{\nu -1}) $
			\item $uf(u) \le (1+|u|^2).$
		\end{enumerate}
If the initial condition $\theta$ belongs to $L_\rho^p(\R^d)$, where $p=\max\{2, \nu\}$, then there exists a solution, $u(t,\cdot) \in L_\rho^p(\R^d))$, in mild form. The solution satisfies for all $q > p$,
	\[
		\sup_{t \in [0,T]} \mathbb{E}\Norm{u(t,\cdot)}^q_{L_\rho^p(\R^d)} \le c_T(1+ \Norm{\theta}^q_{L_\rho^p(\R^d)}),  \quad \forall \; T>0.
	\]
	\end{theorem}

It is a known result that
	\[
		P_t\phi(\theta) = \mathbb{E}\phi(u_\theta(t,\cdot)), \; t\ge 0, \; \phi \in \mathscr{B}_b
	\]
defines a Markovian transition semigroup on the space $\mathscr{B}_b$ of bounded measurable functions. It also satisfies the Feller property and is stochastically continuous. The following result gives conditions for the existence of an invariant measure for $P_t$.

	\begin{theorem}[Theorem 2:]
		Under the conditions of Theorem \ref{T:AMth1}, let $(P_t)_{t\ge0}$ denote the stochastically continuous Feller semigroup in $L^2_\rho(\R^d)$ associated with the solutions $(u_\theta(t,\cdot))_{ t \ge 0)}.$ If $\varrho > \bar{\varrho} + d$ for some $\bar{\varrho} > d$, and if the solution is bounded in probability in $L^{p\nu}_{\bar{\varrho}}(\R^d)$ for a suitable initial condition, $\theta_0$, ie
			\[
				\forall \epsilon >0, \; \exists R>0, \; \forall t>0,\; : \; \mathbb{P}( \{ \Norm{u_{\theta_0}(t,\cdot)}_{L^{p\nu}_{\bar{\varrho} }} \ge R \}) < \epsilon,
			\]
		then there exists an invariant measure for $(P_t)_{t\ge 0}$ on $L^p_{\rho}(\R^d)$.
	\end{theorem}

The following, and last, result from \cite{AM03} gives the standard sufficient condition to verify the bounded in probability condition.

	\begin{theorem}[Theorem 3:]
First, let $p=\max\{2,\nu\}$ and $\varrho >d$. For $p>2$, let $c(p)$ denote the universal Burkholder-Davis-Gundy constant. Assume that
		\begin{enumerate}
			\item $|\sigma(x) - \sigma(y)| \le c_\sigma |x-y|$
			\item $ |f(x)-f(y)| \le c_{\nu}|x-y|(1+|x|^{\nu -1} + |y|^{\nu -1}) $
			\item There exists $c_{f,k}>0$ such that $uf(u) \le c_{f,k} - k|u|^2$, where
				\[
					k > c(p)^4 \frac{c_\sigma^4}{16\lambda}, \quad \text{resp. } k > c(p)^2 \frac{ac_\sigma^2}{2}
				\]
			in the cylindrical resp. nuclear case, is satisfied. Then
				\[
					\sup_{t \ge 0} \mathbb{E}\Norm{u_{\theta_0}(t,\cdot)}^p_{L_\rho^p(\R^d)} < \infty
				\]
			for every bounded and continuous function $\theta_0:\R^d \to \R.$
		\end{enumerate}
	\end{theorem}

\subsection{Existence and Uniqueness of Invariant Measures for Stochastic Reaction-Diffusion Equations in Unbounded Domains. Misiats, Stanzhytskyi and Yip \cite{MSY15}}

This paper focuses on the SPDE
	\[
		\begin{cases}
			\frac{\partial}{\partial t} u(t,x) = A u(t,x) + f(x,u(t,x)) + \sigma(x,u(t,x)) \dot{W}(t,x), & t>0, \; x \in G
			\\ u(0,x) = u_0(x) & x \in \R^d
		\end{cases}
	\]
where $G \subseteq \R^d$ is a possibly unbounded domain and $A$ is an elliptic operator. The noise $W$, is defined to be the same as the Nuclear case in \cite{AM03}.

	\begin{remark}
		The above seems to be for general operator $A$ and domain $G$. However, in section 2 they focus on the case of $A = \Delta$ and $G=\R^d$.
	\end{remark}

Here is the first result verifying the existence of a invariant measure.
	\begin{theorem}[Theorem 1:]\label{T:MSY15th1}
Let $d \ge 3$. Assume that:
		\begin{enumerate}
			\item $\sigma : \R^d \times \R \to \R$ satisfies $|\sigma(x,u_1) - \sigma(x,u_2)| \le c|u_1-u_2|$ and that $|\sigma(x,u)| \le \sigma_0 \in \mathbb{N}$.
			\item $f: \R^d \times \R \to \R$ satisfies $|f(x,u_1) - f(x,u_2)| \le c|u_1-u_2|$ and there exists a $\varphi \in L^1(\R^d) \cap L^\infty(\R^d)$ such that $|f(x,u)| \le \varphi(x)$.
		\end{enumerate}
Let $u(t,x)$ be the solution of the SPDE with $\mathbb{E}\Norm{(u(0,x)}_{L^2(\R^d)}^2 < \infty$. Then we have that
	\[
		\sup_{t \ge 0 } \mathbb{E}\Norm{(u(t,x)}_{L^2_{\rho}(\R^d)}^2
	\]
	\end{theorem}

\begin{remark}
	Theorem \ref{T:MSY15th1} implies that it allows random initial conditions but this is never explicitly stated.
\end{remark}

Here is the second result presented in this paper. However I am unsure if it applies to our case. This result deals with what they refer to as $\textit{nonlocal stoachastic reaction-diffusion equations.}$
	\[
		\begin{cases}
			\frac{\partial}{\partial t} u = Au+ f(u)+ \sigma(u) \dot{W}, & t>0
			\\ u(0,x) = u_0(x) .
		\end{cases}
	\]
In this example they assume that $f : L^2(\R^d) \to L^2(\R^d)$ is a nonlocal map of the form
	\[
		f(u) = \int_{\R^d} g(x,y,u(t,y)) \ud y.
	\]
The function $g$ is not defined in the paper. They mention that a particular $f$ which is used in the model of nonlocal consumption of resources, as well as the nonlocal Fisher-KPP equation, is
	\[
		\begin{cases}
		f(u) = u(1-\Norm{u}) & \text{if} \Norm{u} \le 1
		\\ 0 & \text{otherwise}.
		\end{cases}
	\]
	\begin{theorem}[Theorem 2:]
		Let $d \ge 3$. Assume that:
			\begin{enumerate}
				\item $\forall u,v \in L^2(\R^d),$ $\Norm{f(u)-f(v)}_{L^2(\R^d)} \le C\Norm{u-v}_{L^2(\R^d)}$
				\item $\forall u,v \in L^2(\R^d),$ $\Norm{\sigma(\cdot,u(\cdot))-\sigma(\cdot,v(\cdot))}_{L^2(\R^d)} \le C\Norm{u-v}_{L^2(\R^d)}$
				\item For some $N>0$, $f(u)=0$ if $\Norm{u}_{L^2(\R^d)} \ge N.$
				\item There exists a $\psi(x) \in L^2(\R^d)$ such that $|\sigma(x,u)| \le \psi(x).$
			\end{enumerate}
		Let $u(t,x)$ be the solution to the SPDE with $u(0,x) = u_0(x) \in L^2(\R^d).$ Then
			\[
				\sup_{t\ge 0}\mathbb{E} \Norm{u(t,x)}_{L^2(\R^d)}^2 < \infty.
			\]
	\end{theorem}

\subsection{Ito's formula in infinite dimension and its application to the existence of invariant measures. Misiats, Stanzhytskyi and Yip \cite{MSY16}}

This paper focuses on the SPDE
	\[
		\begin{cases}
			\frac{\partial}{\partial t} u(t,x) = A u(t,x) + f(x,u(t,x)) + \sigma(x,u(t,x)) \dot{W}(t,x), & t>0, \; x \in G
			\\ u(0,x) = u_0(x) & x \in \R^d
		\end{cases}
	\]
where $G \subseteq \R^d$ is a possibly unbounded domain and $A$ is an elliptic operator. The noise $W$, is defined to be the same as the Nuclear case in \cite{AM03} (white in time and colored in space). They mention that they focus on the case $G=\R^d$. Also, in Theorems 1 and 2, they focus on $A =\Delta$ and these are the two theorems we are focusing on anyway.

The first result deals with an analytically weak solution instead of a mild solution.

	\begin{definition}
		An $\mathcal{F}_t$ adapted random process $u(t,\cdot) \in H^1(\R^d) = W^{1,2}(\R^d)$ is called an analytically weak solution of the SPDE if for any $\psi \in H^1(\R^d)$:
			\begin{align*}
				\int_{\R^d} u(t,x)\psi(x) \ud x = & \int_{\R^d} u(0,x)\psi(x) \ud x - \int_0^t \int_{\R^d} \nabla u(s,x) \nabla \psi(x) \ud x \ud s + \int_0^t \int_{\R^d} f(u(s,y)) \psi(x) \ud x
				\\ &+ \sum_{k=1}^\infty \int_0^t \int_{\R^d} \sigma(x,u(s,x)) e_k(x) \sqrt{a_k} \psi(x) \ud x \ud B_k(s)
			\end{align*}
	\end{definition}

Here is the first result.

	\begin{theorem}[Theorem 1:]
For $ t \ge 0$, let $u(t,\cdot) \in H^1(\R^d)$ be an analytically weak solution of the SPDE with initial condition $u(0,x) = u_0(x) \in L^2(\R^d)$. Assume that:
		\begin{enumerate}
			\item $A = \Delta$ and $d \ge 3$
			\item $f(x,u) : \R^d \times \R \to \R$ is such that $|f(x,0)| \in L^2(\R^d)$.
			\item There exists an $L>0$ such that
				\[
					|f(x,u) - f(x,y)|, \; |\sigma(x,u) - \sigma(x,v)| \le L|u-v|
				\]
			\item There exists an $M > 0$ and a nonnegative $\eta(x) \in L^1(\R^d)$ such that
				\[
					uf(x,u) \le \eta(x) \quad \text{for } |u| > M \text{ and } x \in \R^d
				\]
			\item There exists $\sigma_0 >0$ such that $|\sigma(x,u)| \le \sigma_0$.
Then
	\[
		\sup_{t \ge 0} \mathbb{E}\Norm{u(t,x)}_{L^2_{\rho}(\R^d)}^2 \le C < \infty
	\]
for any admissible weight $\rho$.
		\end{enumerate}
	\end{theorem}
The next result deals with the ordinary mild solution. In addition the assumption that
	\[
		\sup_k \sup_{x\in \R^d} \Norm{e_k}_{L^\infty(\R^d)} \le 1
	\]
must be made.

	\begin{theorem}[Theorem 2:]
Let $u(t,x)$ be the mild solution of the SPDE in $L^2_{\rho}(\R^d)$ with initial condition $u(0,x) \ in L^2(\R^d)$. Assume that:
		\begin{enumerate}
			\item $A = \Delta$ and $d \ge 3$
			\item The noise $W$ satisfies $	\sup_k \sup_{x\in \R^d} \Norm{e_k}_{L^\infty(\R^d)} \le 1$
			\item $|f(x,0)| \le \varphi(x)$ and $|\sigma(x,0)| \le \varphi(x)$ where
				\[
					\frac{\varphi(x)}{\sqrt{\rho(x)}} \in L^{\infty(\R^d)} \cap L^2(\R^d)
				\]
			\item $f$ and $\sigma$ satisfy
				\[
					|f(x,u) - f(x,y)|, \; |\sigma(x,u) - \sigma(x,v)| \le L\varphi(x)|u-v|
				\]
		\end{enumerate}
Then for sufficiently small $L$, the solution $u(t,x)$ satisfies
	\[
			\sup_{t \ge 0} \mathbb{E}\Norm{u(t,x)}_{L^2_{\rho}(\R^d)}^2 < \infty
	\]
for any admissible weight.
	\end{theorem}

\section{A comparison of our paper to previous results.}
	\[
		\begin{cases}
			\frac{\partial}{\partial t} u(t,x) = \Delta u(t,x) + f(x,u(t,x)) + b(x,u(t,x)) \dot{W}(t,x), & t>0, \; x \in \R^d
			\\ u(0,x) = u_0(x) & x \in \R^d
		\end{cases}
	\]

First, from a purely educational note, this paper uses Dalang theory to answer a question that has previously only been addressed using DaPrato theory. As far as I know, there has not been another paper that uses the theory of Dalang to show the existence of an invariant measure. Because of this, the noise presented in our paper is constructed in a different way. The noise using the DaPrato approach is constructed using a "basis" of i.i.d. Wiener processes $(B_k(t))_k$. Our noise is constructed by first starting with a positive definite functional
	\[
		J(\varphi, \psi) = \int_0^\infty \ud s \int_{\R^d} \ud x \int_{\R^d} \ud y \; \varphi(s,y) f(x-y) \psi(s,x) \quad \varphi, \psi \in C^\infty_0(\R^{d+1})
	\]
where $\Gamma(\ud x) = f(x) \ud x$ is a non-negative and non-negative definite tempered measure. We then apply a well known result that every positive definite functional is the covariance functional for some zero-mean Gaussian processes. Through the theory of worthy martingale measures presented in \cite{DalangEtc09Minicourse}, we arrive at our noise $W$ which is defined as a particular $L^2$ limit of a sequence of the Gaussian processes defined from $J$ above.

In 1992, Daprato and Zabczyk \cite{DZ92} established a procedure for proving the existence of an invariant measure for those SPDEs whose partial differential operator generates a semigroup which satisfies a particular compactness property. The procedure is as follows:
	\begin{enumerate}
		\item \label{S:step 1} Show the existence of a Markovian solution of the SPDE in a certain function space and verify that its transition semigroup satisfies the Feller property;
		\item \label{S:step 2} Show that the solution to the SPDE is bounded in probability.
	\end{enumerate}
Daprato \textit{et.al.} \cite{DZ96} have verified Step \ref{S:step 1} to be true for a wide variety of differential operators. Step \ref{S:step 2} on the other hand is often nontrivial and requires careful selection of the coefficients of the SPDE. When dealing with the weighted $L^2_{\rho}(\R^d)$ function space, for some \textit{admissible weight} $\rho$, a common trick used to verify Step \ref{S:step 2} is to prove that
	\begin{equation}\label{E:bipc}
		\sup_{t \ge 0} \mathbb{E}(|u(t,\cdot)|_\rho^2) < \infty.
	\end{equation}
This was the case in \cite{TessitoreZabczyk98M,AM03,MSY15,MSY16} and it was also the technique used in this paper as well.

Tessitore \textit{et.al.} \cite{TessitoreZabczyk98M} considered the case where $f\equiv 0$, $b(x,\cdot)$ be Lipschitz continuous and $|\sigma(x,0)| < b_0<\infty$ for each $x\in\R^d$.  In addition, they required the existence of a $p\ge1$ such that $(1-1/p)(d/2) < 1$ and $\gamma \in L^p(\R^d)$ where $\gamma$ is the spectral density of their noise. In Theorem 3.1 ibid. they proved that if there exists two weights, $\rho$ and $\hat{\rho}$, such that $\rho(\hat{\rho})^{-1} \in L^1(\R^d)$ and there exists an element $\hat\varphi \in L^2_{\hat{\rho}}(\R^d)$ such that $u^{\hat{\varphi}}(t,\cdot)$ is bounded in probability in $L^2_{\hat{\rho}}(\R^d)$ then there exists an invariant measure for the SPDE in $L^2_{\rho}$. In Theorem 3.3 ibid. it is verified that for $d \ge 3$ and when the Lipschitz constant of $b$ is sufficiently then for any admissible weight $\rho$, the condition \eqref{E:bipc} is satisfied for the constant one initial condition, $u(0,x)=1$ and thus proving the existence of an invariant measure for this case.

Assing \textit{et.al.} \cite{AM03} also considers the case where $b(x,\cdot)$ is Lipschitz continuous. However in this paper they only consider the weight $\rho(x) = (1+|x|^2)^{-\varrho/2}$. On the the other hand, Assing \textit{et.al.} improves on the results from Tessitore \textit{et.al} \cite{TessitoreZabczyk98M} by allowing all spatial dimensions $d\ge 1$ and by both considering a noise that is white in time and space when $d=1$ and white in time and colored in space for dimensions $d \ge 1$. In addition they allow a non-zero $f$ that satisfies several growth bounds which can be found in Theorem 3 ibid. In this theorem, they prove that for any bounded and continuous initial condition $\theta_0(x) = u(0,x)$ and for the particular choice of $\rho$ given above that
	\[
		\sup_{t \ge 0} \mathbb{E}(|u(t,\cdot)|_{p,\varrho}^p) < \infty
	\]
where
	\[
		|u(t,\cdot)|_{p,\varrho}^p = \int_{\R^d} |u(t,x)|^p (1+|x|^2)^{-\varrho/2} \ud x.
	\]
Hence the existence of an invariant measure is guaranteed under these conditions.

Lastly, Misiats \textit{et.at.} \cite{MSY16} have considered the case of bounded coefficients $f$ and $b$ in the sense that for and admissible weight, there must exist a function $\varphi(x)$ such that $\varphi(x)(\rho(x))^{-1/2} \in L^\infty(\R^d) \cap L^2(\R^d)$ and
	\[
		|f(x,0)| \le \varphi(x) \quad \text{and} \quad |b(x,0)| \le \varphi(x).
	\]
In addition, $f$ and $b$ must satisfy the following Lipschitz type bound growth bound:
	\[
		|f(x,u) - f(x,y)|, \; |b(x,u) - b(x,v)| \le L\varphi(x)|u-v|.
	\]
In this paper, the spatial dimension $d\ge3$ and the noise is considered to be white in time and colored in space. The noise in this paper is defined analogously as the noise in Assing \textit{et.al.} \cite{AM03} for the nuclear case. With the above assumptions, Misiats \textit{et.al.} proved in Theorem 2 ibid. that for a suffieciently small Lipschitz constant $L$ that condition \eqref{E:bipc} is satisfied for any initial condition $u_0(x) = u(0,x) \in L^2(\R^d)$ thus improving the choice of the initial condition when compared to \cite{TessitoreZabczyk98M,AM03}.

We continue this pattern of trade-offs in order to further relax some certain conditions. For starters, we only consider the case where $f\equiv0$ but, as an improvement to \cite{MSY16}, we allow for the same restriction presented on $b$ as in \cite{TessitoreZabczyk98M}. In addition, we allow for all spatial dimensions, as was also the case in \cite{AM03}, but as an improvement to that paper, we allow for any admissible weight. A contribution that we add which is an improvement to all the previously mentioned results is the verification of condition \eqref{E:bipc} for any initial conditions, $u_0(x) = u(0,x)$ such that  $u_0(x) \in L^2_{\rho}(\R^d)$. In addition, we verify condition \eqref{E:bipc} for measure valued initial conditions, the Dirac delta measure included, that satisfy some specific properties. Hence this paper offers conditions which allow us to expand our choices for the initial conditions.


\bigskip
{ \color{magenta}
Here are some suggestions:
\begin{itemize}
	\item Please find more nontrivial admissible initial conditions that are not in $L_\rho^2(\R^d)$.
	\item Extending this work to more general differential operators following \cite{Huang17}. The
		moment formula there extend our formula in \cite{CK15SHE} and allows a Lipschitz drift term.
	\item We should emphasize two contributions: (1) Noise, we can include Riesz kernel, which is not
		a trace or nuclear operator. (2) Initial data, we only require that it is a measure such that
		the solution to the homogeneous equation is well defined for all time and space, namely,
		\begin{align*}
			J_0(t,x):=(\mu * p_t)(x)<\infty, \qquad \text{for all $t>0$ and  $x\in\R^d$}
		\end{align*}
		and
		\begin{equation*}
			J_0(t,\cdot)	\in L_\rho^2(\R^d) \qquad \text{for some $t>0$.}
		\end{equation*}
	\item The above conditions for initial data are general enough, which include
		$L_\rho^2(\R^d)$-valued initial data as a special case. In the proof and statement of our
		results, we may simply focus on this general case and leave the $L_\rho^2(\R^d)$ as an
		intermediate step in the proof.
\end{itemize}
}

\begin{small}
	\addcontentsline{toc}{section}{Bibliography}
	\begin{thebibliography}{999}

		\bibitem{AM03}
		Assing, Sigurd and Ralf Manthey.
		\newblock Invariant measures for stochastic heat equations with unbounded coefficients.
		\newblock {\it Stochastic Process. Appl.} 103 (2003), no. 2, 237--256.

		\bibitem{CarMol94}
		Carmona, Ren\'e~A.  and Stanislav~A. Molchanov.
		\newblock Parabolic Anderson problem and intermittency.
		\newblock {\em Mem. Amer. Math. Soc.}, 108 (1994), no. 518, viii+125 pp.

		\bibitem{Cer01}
		Cerrai, Sandra.
		\newblock {\it Second order PDE's in finite and infinite dimension.}
		\newblock A probabilistic approach. Lecture Notes in Mathematics, 1762. Springer-Verlag, Berlin, 2001. x+330 pp.

		\bibitem{ChenDalang13Heat}
		Chen, Le and Robert C. Dalang.
		\newblock Moments and growth indices for nonlinear stochastic heat equation
		with rough initial conditions.
		\newblock {\em Ann.\ Probab.}  43 (2015), no. 6, 3006--3051.

		\bibitem{CHN16Density}
		Chen, Le, Yaozhong Hu and David Nualart.
		\newblock Regularity and strict positivity of densities for the nonlinear stochastic heat equation.
		% \newblock {\em Preprint at arXiv:1611.03909,} 2016.
		\newblock {\em Mem. Amer. Math. Soc.}, 2019, to appear.

		\bibitem{CH16Comparison}
		Chen, Le and Jingyu Huang.
		\newblock Comparison principle for stochastic heat equation on $\R^d$.
		\newblock {\em Ann.\ Probab.}  47 (2019), no. 2, 989--1035.

		\bibitem{CK15SHE}
		Chen, Le and Kunwoo Kim.
		\newblock Nonlinear stochastic heat equation driven by spatially colored noise: moments and intermittency.
		\newblock {\em Acta Math. Sci. Ser. B}, 39B (2019), No. 3, 645--668.

		\bibitem{DalangEtc09Minicourse} Dalang, Robert, Davar Khoshnevisan, Carl
		Mueller, David Nualart, and Yimin Xiao
		\newblock {\it A minicourse on stochastic partial differential equations}.
		\newblock Held at the University of Utah, Salt Lake City, UT, May 8–19, 2006.
		Edited by  Khoshnevisan and Firas Rassoul-Agha. Lecture Notes in Mathematics,
		1962.
		\newblock Springer-Verlag, Berlin, 2009. xii+216 pp.

		\bibitem{DalQuer}
		Dalang, Robert C. and Llu\'is Quer-Sardanyons.
		\newblock Stochastic integrals for spde's: a comparison.
		\newblock {\it Expo. Math.} 29 (2011), no. 1, 67--109.

		\bibitem{DZ92}
		Da Prato, Giuseppe and Jerzy Zabczyk.
		\newblock {\it Stochastic equations in infinite dimensions.}
		Encyclopedia of Mathematics and its Applications, 44. Cambridge University Press, Cambridge, 1992. xviii+454 pp.

		\bibitem{DZp92}
		Da Prato, Giuseppe and Jerzy Zabczyk.
		\newblock {\it Invariant measures for semilinear stochastic equations.}
		Stochastic Anal. Appl. 10, 387-408.


		\bibitem{DZ96}
		Da Prato, Giuseppe and Jerzy Zabczyk.
		\newblock {\it Ergodicity for infinite-dimensional systems. }
		\newblock London Mathematical Society Lecture Note Series, 229. Cambridge University Press, Cambridge, 1996.

		\bibitem{Huang17}
		Huang, Jingyu
		\newblock {\it On stochastic heat equation with measure initial data}
		\newblock {\it Electron. Commun. Probab.} 22 (2017) . No. 40, 1-6.

		\bibitem{MS99}
		Maslowski, Seidler.
		\newblock {\it On sequentially weakly Feller solutions to SPDEs.}
		\newblock Rend. Lincei Mat. Appl. 10, 69-78

		\bibitem{MSY16}
		Misiats, Oleksandr, Oleksandr Stanzhytskyi, and Nung-Kwan Yip.
		\newblock Existence and uniqueness of invariant measures for stochastic reaction-diffusion equations in unbounded domains.
		\newblock {\it J. Theoret. Probab.} 29 (2016), no. 3, 996--1026.


		\bibitem{MSY15}
		Misiats, Oleksandr, Oleksandr Stanzhytskyi, and Nung-Kwan Yip.
		\newblock Ito's Formula in infinite dimension and its application to the existence of invariant measures.

		\bibitem{NIST2010}
		F.~W.~J. Olver, D.~W. Lozier, R.~F. Boisvert, and C.~W. Clark, editors.
		\newblock {\em N{IST} handbook of mathematical functions}.
		\newblock U.S. Department of Commerce National Institute of Standards  and Technology, Washington, DC. Cambridge Univ. Press, Cambridge, 2010.

		\bibitem{TessitoreZabczyk98M}
		Tessitore, Gianmario and Jerzy Zabczyk.
		\newblock Invariant measures for stochastic heat equations.
		\newblock {\it Probab. Math. Statist.} 18 (1998), no. 2, Acta Univ. Wratislav. No. 2111, 271--287.

		\bibitem{Walsh} Walsh, John B.
		\newblock {\it An Introduction to Stochastic Partial Differential Equations}.
		\newblock In: \`Ecole d'\`et\'e de probabilit\'es de
		Saint-Flour, XIV---1984, 265--439.
		Lecture Notes in Math.\ 1180, Springer, Berlin, 1986.

	\end{thebibliography}
\end{small}

\end{document}
